\chapter{Atmospheric Phenomena and Climate-Model Data}
\label{ch:atmospheric}

The previous chapters isolate the numerical core of the thesis: randomized
approximation, mixed precision, least-squares geometry, and regularized
covariance transport. Before returning to the concrete bias-correction workflow,
it is useful to make explicit what kind of atmospheric structure the application
is meant to represent. The purpose of this chapter is therefore contextual. It
does not shift the thesis away from numerical linear algebra; rather, it explains
why the covariance and dependence operators used later are meaningful objects in
climate-model post-processing.

Atmospheric data are multivariate by nature. Temperature, precipitation, wind,
humidity, radiation, pressure, and derived extremes are not independent columns in
a table. They are coupled through thermodynamic, radiative, hydrological, and
dynamical processes. A statistical correction method that adjusts only marginal
distributions may therefore improve each variable separately while leaving the
joint state physically implausible. This is the basic reason why multivariate
bias correction is a relevant application for the covariance-transport methods
studied in this thesis.

\section{Weather, Climate, and Extremes}

The climate-science motivation is strongest for variables and events whose
impacts depend on joint behaviour. The IPCC Sixth Assessment Report emphasizes
that human-induced climate change is already affecting weather and climate
extremes, including hot extremes, heavy precipitation, droughts, storms, and
compound events \cite{ipcc2021ar6wg1}. These phenomena are not only changes in
individual means. They involve changes in tails, spatial extent, persistence,
co-occurrence, and physical drivers.

Temperature extremes provide the simplest example. A bias correction that
matches daily temperature quantiles may still distort the relationship between
daily minimum and maximum temperature, humidity, radiation, or wind. Heavy
precipitation is similar. Its statistical behaviour depends not only on the
marginal distribution of rainfall, but also on circulation, moisture availability,
convective activity, orographic forcing, and temporal clustering. Droughts add a
further layer: precipitation deficits, evapotranspiration, soil moisture, and
radiative forcing interact over longer time scales. Compound events, such as
concurrent heat and drought, make the multivariate nature of the problem
explicit.

For the present thesis, this motivates a narrow but useful modelling stance. We
do not attempt to model all atmospheric dynamics. Instead, we ask whether a
dependence-correction layer based on covariance operators can preserve or improve
some second-order structure after marginal correction. This is less ambitious
than a full physical or distributional model, but it gives a concrete numerical
object whose approximation and finite-precision behaviour can be studied.

\section{Regional Reanalysis and Central Europe}

The new local literature also gives a more concrete regional setting. Beranova
et al. validate an ALADIN-Climate regional reanalysis over Europe, with emphasis
on Central Europe \cite{beranova2025aladin}. The model is based on the ALARO
configuration, uses a convection-permitting horizontal resolution of \(2.3\) km,
and is evaluated over the period 1990--2014 against ground-based meteorological
observations. The variables considered are exactly the kinds of fields that make
multivariate post-processing nontrivial: daily temperature, precipitation,
relative humidity, wind speed, and global radiation.

The validation results are relevant to the thesis for two reasons. First, they
show that high-resolution regional reanalyses can be valuable data sources for
Central European climate studies, especially where global reanalyses are too
coarse to resolve local terrain and mesoscale structure. Second, they show that
even high-resolution reanalyses have systematic, variable-dependent biases. In
the ALADIN-Climate validation, mean daily temperature biases are relatively small,
but minimum daily temperature is generally overestimated and maximum daily
temperature underestimated. Winter precipitation totals are overestimated by
nearly \(50\%\) on average, partly through an overestimation of wet days. Relative
humidity, wind speed, and global radiation also show systematic seasonal and
spatial biases \cite{beranova2025aladin}.

These findings are useful because they give the application chapter a realistic
interpretation of the columns in the data matrix. The variables are not abstract
features. They are atmospheric quantities with different error mechanisms:
microphysics and wet-day frequency for precipitation, boundary-layer and
roughness effects for wind, cloud and radiation processes for global radiation,
and diurnal-cycle errors for minimum and maximum temperature. A multivariate
correction method should therefore be evaluated not only by numerical error, but
also by whether it preserves relationships among these physically linked
quantities.

\section{Variables as a Multivariate State}

The data model used later can now be read physically. The application chapter
uses historical model data, future model data, and historical reference data,
\[
	X_h \in \R^{n \times p}, \qquad
	X_f \in \R^{m \times p}, \qquad
	Y_h \in \R^{n \times p}.
\]
The column dimension \(p\) is not fixed by the theory. It may encode variables at
one site, several nearby locations, lagged time features, seasonal indicators, or
blocks of spatial grid points. The same algebraic object therefore covers several
application designs:
\begin{itemize}
	\item intervariable dependence, such as temperature--humidity or
	      precipitation--radiation relationships;
	\item spatial dependence, such as correlation between nearby grid cells or
	      stations;
	\item temporal dependence, represented through lagged columns;
	\item local multivariate states combining variables, locations, and lags.
\end{itemize}

After marginal correction, the randomized covariance-transport method acts on a
latent rank-normal representation. The covariance operators
\[
	C_X = \frac1n Z_X^\top Z_X,
	\qquad
	C_Y = \frac1n Z_Y^\top Z_Y
\]
should therefore be interpreted as summaries of dependence in this latent state
space. They are not a full physical description of the atmosphere. They are a
numerically tractable representation of second-order structure after the
univariate part of the bias correction has already been handled.

This interpretation also clarifies why the spectrum of the covariance operator
matters. Atmospheric fields often contain coherent spatial and intervariable
patterns, but they also contain localized extremes and small-scale variability.
If the dominant dependence is well approximated by a low-dimensional subspace,
then randomized low-rank methods may be effective. If the spectrum is flat or the
important structure lies in weak but physically meaningful modes, low-rank
compression can fail even when the randomized algorithm behaves as expected
mathematically.

\section{Implications for Covariance Transport}

The atmospheric context gives the covariance-transport map a clear but limited
role. The map
\[
	T_\lambda
	=
	(C_X+\lambda I)^{-1/2}(C_Y+\lambda I)^{1/2}
\]
is a dependence-correction operator in latent space. It is designed to move the
second-order structure of the corrected model sample toward that of the
reference sample. It can therefore plausibly affect intervariable correlations,
spatial covariance, and some low-order signatures of compound behaviour.

At the same time, its limitations are structural. Covariance transport does not
model causality, atmospheric circulation, event sequencing, or nonlinear tail
dependence. It does not by itself guarantee that a corrected heatwave, drought,
or heavy-precipitation episode is dynamically plausible. This is why the thesis
uses covariance transport as a numerical core rather than as a complete
climate-science model. The method is valuable only if its narrower target is
stated honestly: a compressed and analyzable dependence-correction layer after
marginal adjustment.

This also explains the diagnostics used later. Operator-level errors such as
\(\|C-\widehat C\|_2\) and
\(\|T_\lambda-\widehat T_\lambda\|_2\) measure whether the numerical method
approximates its own target. Application-facing diagnostics, such as correlation
error, spatial or temporal autocorrelation error, and multivariate discrepancy
measures, ask whether that target is useful for the atmospheric data. Both levels
are necessary. A small transport error is not automatically a successful climate
correction, and a useful climate correction may tolerate an operator error that
looks large in spectral norm.

\section{Role in the Thesis}

The role of this chapter is to bridge the mathematical and application-facing
parts of the thesis. The central contributions remain the perturbation analysis,
the randomized low-rank approximation, and the mixed-precision error
decomposition. Atmospheric phenomena enter as the reason why the data matrices
have multivariate structure and why dependence correction is worth studying.

In the next chapter, this context is specialized to a bias-correction workflow.
The atmospheric variables become columns of the data matrices, marginal
adjustment is separated from dependence correction, and the covariance-transport
operator from \Cref{ch:randcovtransport} becomes the numerical layer tested in
the climate-model application.
